\section{Exact scaling of the canonical solution}
\label{S:exact-canonical-scaling}

We first rewrite the second moment at a large time $T$ as a second moment at
time one, for an equation with the small parameter $\delta = T^{-\kappa}$. The
proof of Proposition~\ref{P:raw-block-point-comparison} uses this for the final
equality in~\eqref{E:raw-block-original-comparison}.

\LeanMarker{L:exact-canonical-scaling}
\begin{lemma}[Canonical law scaling]\label{L:exact-canonical-scaling}
  Under Hypothesis~\ref{H:homog}, put $a = \alpha_0$, $\kappa =
  (4-\alpha-a)/(3-\alpha)$, and, for $T > 0$,
  \[
    \delta=T^{-\kappa},\qquad
    r_T=T\delta=T^{(a-1)/(3-\alpha)}.
  \]
  The law of $u(Tt,r_Tx)$ is the canonical expansion solution $V_\delta(t,x)$
  with data $(1,0)$ of
  \begin{equation}\label{E:exact-fast-equation}
    \delta^2\partial_t^2V_\delta
    =\Delta V_\delta+\sqrt\delta\,\W V_\delta,
  \end{equation}
  where $\W$ has the original covariance. In particular,
  \begin{equation}\label{E:exact-moment-scaling}
    T^{-\kappa}\log\E|u(T,0)|^2
    =\delta\log\E|V_\delta(1,0)|^2.
  \end{equation}
  The assertion is equality in law; it assumes no pathwise uniqueness for a full
  rough solution.
\end{lemma}

\begin{proof}
  We first scale the regularized noise and the Green measure. We check that
  the coefficient at each ordered vertex is the one in the fast equation. We
  then pass the finite-order law identities to the canonical series at fixed
  dilation.


  \medskip
  \noindent\textit{Step 1.~}We scale the source covariance and the Green
  measure.
  At positive widths, time and space dilation of the Gaussian source gives
  covariance amplitude $T^{-a}r_T^{-\alpha}$. This follows distributionally from
  the homogeneous Radon covariance, including spatial white noise. Its source
  cutoffs transform by $K_\eta(Ts) = T^{-a}K_{\eta/T}(s)$ and
  $\epsilon\mapsto\epsilon/r_T^2$; the joint maximum-time and summed-space width
  laws are preserved at every vertex. The Green measure has the scaling
  \[
    G_d(T\tau,r_T\ud z)
    =r_TG_d(\tau/\delta,\ud z).
  \]
  The mass on the left is $T\tau$. The mass on the right is $r_T\tau/\delta$,
  which is the same because $r_T=T\delta$. The spatial support also matches:
  after dividing by $r_T$, the original propagation radius is
  $T\tau/r_T=\tau/\delta$.
  It follows from its degree $1-d$ distributional scaling, also for the sphere
  measure in dimension three. Each terminal Duhamel vertex thus contributes
  $Tr_T\,T^{-a/2}r_T^{-\alpha/2}$. Substitution of the displayed positive real
  powers gives
  \[
    T^{-a/2}r_T^{2-\alpha/2}=\sqrt\delta,\qquad
    T^{1-a/2}r_T^{1-\alpha/2}=\delta^{-1/2}.
  \]

  \medskip
  \noindent\textit{Step 2.~}We match the scaled vertex with the fast-equation
  vertex.
  To check the powers, substitute $r_T=T\delta$ into the second vertex factor.
  It becomes $T^{2-(a+\alpha)/2}\delta^{1-\alpha/2}$. The identity
  $(3-\alpha)\kappa=4-\alpha-a$ and $\delta=T^{-\kappa}$ reduce this to
  $\delta^{-1/2}$. Multiplying this identity by $\delta$ gives the first
  equality in the preceding display.
  The Green kernel for the forcing in~\eqref{E:exact-fast-equation} is
  $\delta^{-1}G_d((t-s)/\delta)$; after multiplication by $\sqrt\delta$ it gives
  exactly that vertex factor. Therefore, every finite list of fixed-order
  coefficients has the same joint law under the dilation, including all
  self-pairings, unequal orders and empty chains. At fixed $T,r_T$, remove Green
  smoothing first in dimension three, with its width divided by $r_T^2$ on the
  scaled side. Lemma~\ref{L:joint-fixed-order} and the absolute $L^2$
  summability in Lemma~\ref{L:vertex-crude} pass the identities to the canonical
  series and its second moment. No large-time and source limit are interchanged.
  For each fixed dilation, both rescaled noise widths still tend to zero, and
  the common finite-order identities are taken before summing. Absolute $L^2$
  summability controls the tails of the series. Finally, evaluate the law
  identity at $t=1,x=0$ and use $T^{-\kappa}=\delta$ to obtain the logarithmic
  moment identity.
  This proves the law assertion and~\eqref{E:exact-moment-scaling}. This proves
  Lemma~\ref{L:exact-canonical-scaling}.
\end{proof}
